\chapter{שימוש בכלים וימד ReAct}

\section{יסודות השימוש בכלים}

מנגנון קשב עצמי, כפי שהוצג בארכיטקטורת ה-\textenglish{Transformer}
\cite{vaswani2017attention}, מהווה בסיס עמוק לעיבוד מידע בקנה מידה גדול.
עם זאת, מוגבלותיו של מודל שפה מבוסס טקסט הופכות למובהקות כאשר משימה
דורשת אינטראקציה עם סביבה חיצונית --- חישובים מתמטיים, שאילתות בסיס
נתונים, או גישה לידע חיצוני עדכני. יכולת השימוש בכלים חיצוניים מהווה
הרחבה טבעית של יכולות הדגם לעבור פעולה חלוקה.

היסוד הראשון של שימוש בכלים הוא זיהוי הצורך בכלי חיצוני. מודל שפה יכול
ללמוד, דרך אימון או הנחייה, להכיר מתי משימה מסוימת עדיפה להיות מחוסנת
בעזרת כלי ספציפי. מחקר בתחום זה הדגיש כי מודלים גדולים מסוגלים לעצמאות
בחירה זו \cite{brown2020language}, במיוחד כאשר הם מוכשרים דרך דוגמאות
דקות (\textenglish{few-shot examples}) המדגימות תבניות של שימוש בכלים.

היסוד השני הוא תרגום כוונה לביצוע כלי. לאחר שמודל השפה קובע כי כלי
מסוים נדרש, עליו לנסח קריאה פונקציונלית (\textenglish{function call})
תקינה --- כלומר, לציין את שם הכלי, פרמטרים ערכים בסדר נכון. הדיוק בשלב
זה הוא קריטי; שגיאה בפרמטר תוביל לכישלון או לתוצאה חסרת משמעות.

היסוד השלישי הוא שילוב התוצאה חזרה לתהליך הנמוכה. כאשר כלי חיצוני
מחזיר תוצאה, המודל חייב לאמת אותה, לפרש את משמעותה, ולשלבה בהמשך
ההנמקה שלו. משלושת יסודות אלו משתרבבת קושיות יישום: בחירה כלים עצמית
דורשת מודל למידה על התפקידים היחסיים של כלים שונים; קריאות כלים תקינות
דורשות בקרת תחביר או כיול זקוק; שילוב חזרה דורש הבנה זכוכה של מבנה
הנתונים המוחזר.

\section{ReAct: חשיבה וביצוע}

מסגרת ה-\textenglish{ReAct} (\textenglish{Reasoning + Acting}) הוצגה כדי
לתת מבנה עקבי לתהליך של חשיבה מדורגת ובצוע כלים באופן תיאום
\cite{yao2023react}. בתוך ה-\textenglish{ReAct}, סוכן פועל בלולאה סגורה:
אנו קוראים ומחשבים, אנו פועלים (קוראים כלי), אנו תצפית בתוצאה, וחוזרים
להחשבה הבאה. מבנה לולאה זה מחקה את ההנמקה האנושית, בה חשיבה והפעלה
אינן נפרדות, אלא משולבות כדי לפתור משימות מורכבות.

הקינמטיקה של \textenglish{ReAct} משתמשת בתבנית הנחיה אחידה שמביעה
במפורש את הנימוק, את ההחלטה בבחירת כלי, את קריאת הכלי, את תוצאת
הכלי, והמשך הנימוק. גישה זו עשתה שתי תרומות משמעותיות: ראשית, היא
יצרה חלון ברור לתוך תהליך ההנמקה של המודל, מה שאפשר למחקרים להעריך
את איכות השלבים הביניים; שנית, היא הוכיחה כי ההנחיה המפורשת של שלבים
אלו משפרת משמעותית את דיוק הפתרון.

בהשוואה לשיטות קודמות בהן מודל היה מתבקש פשוט להוציא תשובה סופית,
\textenglish{ReAct} מסירה את הנטל מן המודל לחשוב כל דבר בו-זמנית. במקום
זאת, המודל מוודא ממצא אחד בכל שלב, רואה את התוצאה, ויכול להתאים את
התוכנית שלו. זה מזכיר תהליכים אנושיים: כשאני משימה מתמטית, אני לא
משחק את כל השלבים בראשי; אני כותב שלב, בודק אותו, ושוקל את התוצאה
לפני שהעבור הבא.

כוחה של \textenglish{ReAct} טמון גם בגמישותה. הרוטו החיצוני (מיהו המודל,
מהם הכלים הזמינים) ניתן לשינוי ללא שינוי המאומנתון --- הן לכך שמודל
אחר יכול לחליף את ה-\textenglish{ReAct} בקלות.

\section{Toolformer וטכניקות בחירת כלים עצמיות}

בעוד ש-\textenglish{ReAct} מוגדר כתבנית הנחיה שימוש בכלים, שאלה מרכזית
נותרה: כיצד מודל שפה לומד \emph{אילו} כלים להשתמש ובאילו נקודות?
\textenglish{Toolformer} \cite{schick2023toolformer} הציע גישה בה המודל
עצמו משתמש בשפה כדי לטפל את הידע שלו על כלים. בשיטה זו, כלים אינם
מוגדרים מראש כמערכת חיצונית קשיחה, אלא החל כחלק מהטקסט עצמו הנוצר
על ידי המודל.

הרעיון הוא זה: אם אנו מהנדסים כלי ודוגמות אימון המדגימות שימוש בו,
המודל יכול ללמוד לחזות קריאות כלים בעבור פולט תוכן בדומה לאופן בו הוא
חוזה מילות טקסט רגילות. זה המסה של ממודול עצמוני: המודל מפתח יכולת
לשפוט כי במצב X, קריאה לכלי Y עם ארגומנט Z היא פרטית בסדר הרחוק שלו.

הטכניקה של בחירת כלים עצמיות מוסיפה שכבה נוספת של מינייטו. לא רק
צריך המודל ללמוד להשתמש בכלי כאשר הוא נקרא בצורה מפורשת; עליו גם
ללמוד לזהות באופן עצמוני מתי השימוש בכלי יהיה חיוני. במחקר
\textenglish{Toolformer}, זה הושג באמצעות פרטים בעלי חיובים חיוביים:
המודל מעריך למימד פנימי כל זמן מה ערך מצופה של קריאה לכלי נתון, וגם
מזניח קריאות בעלות ערך שלילי או נמוך.

מגבלה משמעותית של גישה זו היא שהיא דורשת אימון מחדש או עדכון משמעותי
של המודל עצמו. זה לא ניתן להחיל בקלות על מודלים קפואים או כמ-אפליקציות
רק-קריאה. עם זאת, הוכחתה של הקונספט גרמה לעבודות עוקבות בתחום בחירת
כלים מתאם.

\section{Chain-of-Thought כהתחלה למחשבה מסודרת}

לפני שמודלי שפה גדולים למדו להשתמש בכלים חיצוניים, הם כבר הוכיחו כי
יכולת לפרק משימות מורכבות לשלבים ביניים משפרת את הדיוק. טכניקה זו,
המוכרת ככלל \textenglish{Chain-of-Thought} (\textenglish{CoT})
\cite{wei2022chain}, דגמה כי פשוט בבקשה ממודל לחשבות את שלבי הדרך שלו
בטרם נתן תשובה סופית גרם לשיפור משמעותי בביצוע במשימות המחייבות
הנמקה מתמטית או לוגית.

ה-\textenglish{Chain-of-Thought} מהווה קדמתון חיוני לשימוש בכלים. בעוד
שגישה זו לא כרוכה בקריאות כלים חיצוניות, היא החזיקה בנפרדות של
הדיוקוריה לשם זה: היא גרמה למודלים להתנהג בדרך שנראית מורגלת בצורה
הוגנת. מודל המוטה להראות את השלבים ביניים שלו באופן סדור יוצר נקודת
התחלה טבעית לאינטגרציה של קריאות כלים אמיתיות לתוך זרימת החשיבה הזו.

בהיקפה הרחב, \textenglish{CoT} היא שיטת הנחיה שבה אנו מבקשים מהמודל
להודיע במפורש את תהליך הגיגו. זה שונה מלחלוטין מגישת ``מלל-אל-מלל''
בה המודל מובקש לפשוט לעבור מקלט לפלט ישירה. כאשר \textenglish{CoT}
משולבת עם שימוש בכלים, אנו משיגים מסגרת בה חשיבה וביצוע משולבים
בשיטה מובנית וחזרוגית.

בהקשר של מודלים גדולים כמו \textenglish{GPT-3} \cite{brown2020language},
יכולת \textenglish{CoT} להופיע באופן טבעי דרך הנחיה לא-מובנית
(unstructured prompting) הוכיחה כי מודלים אלו כבר חזקו בדרכם להנמקה
מסדר-גבוה. הציבור על \textenglish{CoT} כטכניקת בחירה לקידום יכולות
השכל היה סיום חברויות נוסף לכיוון של שימוש בכלים יותר מובנה.

נושא קריטי שלא נמתך במוגבלויות זו הוא הטכנולוגיה של תחבור ההנחיות
(\textenglish{prompt engineering}). עבודה עתידית בתחום, כולל תקנים כמו
\textenglish{Model Context Protocol} \cite{anthropic2024mcp}, מתמקדת
בסטנדרטיזציה של האופן בו מודלים מתקשרים עם כלים חיצוניים בצורה
מובנית וחוזרת.

\begin{figure}[htbp]
    \centering
    \includegraphics[width=0.85\textwidth]{../assets/attention_complexity.png}
    \caption{השוואת מורכבות חישובית: תשומת לב סטנדרטית \textenglish{$O(n^{2})$}, תשומת לב לינארית \textenglish{$O(n \log n)$}, ורקורנטי \textenglish{$O(n)$}.}
    \label{fig:attention_complexity}
\end{figure}
